% Friction lookup table v2 (baseline C2) -- Slide-to-Target, MIT 6.4210
% Build the row files first:  python3 lookup_table.py   (writes rows_*.tex and friction_lookup.csv)
% Then:                       pdflatex friction_lookup_table.tex  (twice, for the table numbers)
%
% DRAFT STATUS (team, please read before citing):
%  * The "Difficulties" list (Section 7) was written fresh for v2 and is TO BE MERGED WITH v1 by the team.
%  * Every row has verified = n: nobody has yet checked a value against its source PDF (Section 8).
%  * The Section 10 decisions were settled by the team on 5 Nov 2026 and are applied here: tier factors
%    2/2.5/3; L06 and Y02 share the pooled 0.25 (f = 1.8); V03 uses the "perpendicular" 0.33; Y04 is
%    dropped from the simulated set. See Section 9 and the comment block in lookup_table.py.
\documentclass[11pt]{article}
\usepackage[utf8]{inputenc}
\usepackage[T1]{fontenc}
\usepackage[a4paper,margin=1.7cm]{geometry}
\usepackage{array}
\usepackage{longtable}
\usepackage{amsmath}
\PassOptionsToPackage{hyphens}{url}
\usepackage[hidelinks]{hyperref}
\setlength{\tabcolsep}{3pt}
\setlength{\LTcapwidth}{\textwidth}
\renewcommand{\arraystretch}{1.08}

\title{Friction lookup table v2 (baseline C2)}
\author{Slide-to-Target team, MIT 6.4210}
\date{\today\ --- draft; Difficulties list to be merged with v1; all rows \texttt{verified = n}}

\begin{document}
\maketitle

\noindent\fbox{\parbox{\dimexpr\textwidth-2\fboxsep-2\fboxrule}{\small
\textbf{Header note.} This document is generated in part by \texttt{lookup\_table.py}: Tables~\ref{tab:main}--\ref{tab:excluded} and the reference list are \texttt{\textbackslash input} from \texttt{rows\_*.tex}, so the numbers here are the numbers in \texttt{friction\_lookup.csv}. The Difficulties list (Section~\ref{sec:difficulties}) is new in v2 and is to be merged with the v1 list by the team. No value has been verified against its source PDF yet (\texttt{verified = n} on every row; checklist in Section~\ref{sec:verify}). The Section~10 decisions of the implementation spec were settled by the team on 5 Nov 2026 and are applied (Section~\ref{sec:placeholders}).}}

\section{What changed in v2}
Lookup v2 covers 17 of the 27 rig pairs (M01--M27), up from 2 in v1, using sources independent of Yu et al. It also gives Yu's three simulated pairs (Y01--Y03, formerly L21--L23; Y04, polyurethane, was dropped from the simulated set on 5 Nov, decision 6) independent lookup values, which miss Yu's measurements by about $2\times$ on two of three, so baseline C2 is no longer correct by construction there (Section~\ref{sec:yu}). No rig cell is an exact kinetic match, and felt, acrylic and glazed tile have no kinetic source at all; that gap is itself a finding. The v1 class-fallback table (old Table~4) is dropped: tiered entries replace it, and a pair with no source falls back to the default prior C0.

\section{Rules}\label{sec:rules}
\begin{enumerate}
 \item \textbf{Kinetic values only.} A static value never becomes an entry. Where a close static value exists it is recorded in \texttt{static\_upper}, since kinetic friction is usually lower.
 \item \textbf{Evidence tiers.} A = exact material pair. B = same object material on a similar surface class. C = different counterface class, or reversed geometry.
 \item \textbf{Spread by tier.} The 90\% interval is $\mu/f$ to $\mu f$, with $f = 2$ (A), $2.5$ (B), $3$ (C) (team decision 3, 5 Nov; the v2 draft had 1.5 / 2 / 2.5); $s_{\text{tab}} = \ln f / 1.645 = 0.421,\ 0.557,\ 0.668$. The lookup value enters the belief as one Gaussian observation of $\ln\mu$ on top of the population prior (\texttt{belief.loglik\_lookup}). Spread B rests on how far apart speed-matched PLA-on-steel studies are (0.23 to 0.49). These factors are a modelling choice, not a measurement. Rows L01--L20 keep their v1 spreads ($f = 1.5$, wider where sources disagree); their tier letter is informational. L06 and Y02 share the pooled acetal-on-steel entry with $f = 1.8$ (decision 5).
 \item \textbf{No Yu et al.} in any lookup row or source list, and nothing derived from their dataset. Their values appear in this document only in Table~\ref{tab:yu}, read from the simulator-truth file.
 \item \textbf{Cross-checks only.} The Engineering ToolBox and other uncited tables never set a value.
 \item \textbf{No gap-filling by asking an LLM.} That would turn C2 into a second LLM condition.
 \item \textbf{Tier ``none''.} A pair with no source gets tier none, blank value columns, and C2 falls back to the default prior C0.
\end{enumerate}
Columns of \texttt{friction\_lookup.csv}: the eleven v1 columns (\texttt{pair\_id} \ldots\ \texttt{notes}) unchanged, then \texttt{role} (secondary for L01--L20, yu\_pair, rig), \texttt{tier} (A, B, C, none), \texttt{static\_upper}, \texttt{source\_speed\_mps} (blank if the source does not state a speed), \texttt{source\_method}, \texttt{verified}. Source numbers in the tables refer to Section~\ref{sec:refs}, which also gives the key used in the CSV.

\section{Table 1: independent values (L01--L20 and Y01--Y03)}
Rows L01--L20 come from the handbook table and are unchanged from v1 except L06, which gains source [15], loses the Yu figure and takes the pooled 0.25 shared with Y02 (decision 5). Rows Y01--Y03 carry the pair IDs of the simulator-truth file and show ``(L21)'' etc.\ for the v1 ID; their values are independent of Yu et al. Y04 (polyurethane) was dropped from the simulated set on 5 Nov (decision 6); its candidate value is in Table~\ref{tab:excluded}.

{\footnotesize
\begin{longtable}{l p{2.8cm} p{2.2cm} r l c c c l l}
\caption{Lookup entries: handbook rows L01--L20 and the four Yu pairs with independent values. $\mu_k$ is the nominal kinetic coefficient, $s$ the log spread, the interval the 90\% range. Flags: ``near 0.50 cap'' = at the top of the simulated range; ``$n$ sources'' = nominal chosen among disagreeing sources.}\label{tab:main}\\
\hline
ID & Object (moving) & Surface & $\mu_k$ & Source range & 90\% interval & $s$ & Tier & Src & Flag\\
\hline
\endfirsthead
\multicolumn{10}{l}{\emph{Table~\ref{tab:main}, continued}}\\
\hline
ID & Object (moving) & Surface & $\mu_k$ & Source range & 90\% interval & $s$ & Tier & Src & Flag\\
\hline
\endhead
\hline
\endfoot
L01 & hard steel block & mild steel plate & 0.42 & single value & 0.28--0.63 & 0.25 & A & [1,2,3] & \\
L02 & aluminium block & mild steel plate & 0.47 & single value & 0.31--0.71 & 0.25 & A & [1,2,3] & near 0.50 cap\\
L03 & copper block & mild steel plate & 0.36 & single value & 0.24--0.54 & 0.25 & A & [1,2,3] & \\
L04 & brass block & mild steel plate & 0.44 & single value & 0.29--0.66 & 0.25 & A & [1,2,3] & \\
L05 & phenolic laminate block & mild steel plate & 0.35 & single value & 0.23--0.53 & 0.25 & A & [1,2] & \\
L06 & acetal (POM) block & mild steel plate & 0.25 & 0.15--0.45 & 0.14--0.45 & 0.36 & A & [4,5,7,15] & 4 sources, pooled with Y02\\
L07 & nylon 6/6 block & mild steel plate & 0.27 & 0.17--0.57 & 0.15--0.49 & 0.36 & A & [6,5,7] & 3 sources\\
L08 & UHMW-PE block & mild steel plate & 0.14 & 0.09--0.25 & 0.088--0.22 & 0.29 & A & [5,6] & \\
L09 & zinc block & cast iron plate & 0.21 & single value & 0.14--0.32 & 0.25 & A & [1,2,3] & \\
L10 & copper block & cast iron plate & 0.29 & single value & 0.19--0.44 & 0.25 & A & [1,2,3] & \\
L11 & brass block & cast iron plate & 0.30 & single value & 0.20--0.45 & 0.25 & A & [1,2,3] & \\
L12 & magnesium block & cast iron plate & 0.25 & single value & 0.17--0.38 & 0.25 & A & [1,2,3] & \\
L13 & bronze block & cast iron plate & 0.22 & single value & 0.15--0.33 & 0.25 & A & [1,2,3] & \\
L14 & mild steel block & cast iron plate & 0.23 & single value & 0.15--0.35 & 0.25 & A & [1,2,3] & \\
L15 & cast iron block & cast iron plate & 0.15 & single value & 0.10--0.23 & 0.25 & A & [1,2,3] & static 1.10\\
L16 & oak block & cast iron plate & 0.49 & single value & 0.33--0.74 & 0.25 & C & [1,2,3] & near 0.50 cap\\
L17 & oak block, grain parallel & oak board & 0.48 & single value & 0.32--0.72 & 0.25 & A & [1,2,3] & \\
L18 & oak block, grain perpendicular & oak board & 0.32 & single value & 0.21--0.48 & 0.25 & A & [1,2,3,8] & \\
L19 & glass block & glass plate & 0.40 & single value & 0.27--0.60 & 0.25 & A & [1,2,3] & \\
L20 & graphite block & glass plate & 0.18 & single value & 0.12--0.27 & 0.25 & A & [1,2] & \\
Y01 (L21) & stainless steel block & ABS sheet & 0.24 & single value & 0.08--0.72 & 0.67 & C & [12] & replaces Yu value\\
Y02 (L22) & stainless steel block & acetal (Delrin) sheet & 0.25 & 0.15--0.45 & 0.14--0.45 & 0.36 & A & [4,5,7,15] & replaces Yu value; pooled with L06\\
Y03 (L23) & stainless steel block & plywood sheet & 0.22 & single value & 0.073--0.66 & 0.67 & C & [18,19] & replaces Yu value\\

\end{longtable}
}

\section{Table 2: static-only, unlabelled and near-miss values}
These are not entries (rule 1). Where a static value is close to a rig pair it is carried in that row's \texttt{static\_upper} as an upper bound. The last block is the static part of the v1 excluded list.

{\footnotesize
\begin{longtable}{p{4.4cm} l l l p{6.2cm}}
\caption{Values found but not usable as kinetic entries: static (tilt tests, lashing-guide design factors, secondhand static tables), unlabelled (static/kinetic not stated), or kinetic on the wrong surface.}\label{tab:static}\\
\hline
Pair & Value & Type & Src & Note\\
\hline
\endfirsthead
\multicolumn{5}{l}{\emph{Table~\ref{tab:static}, continued}}\\
\hline
Pair & Value & Type & Src & Note\\
\hline
\endhead
\hline
\endfoot
furniture felt on melamine-faced base & 0.41 & static & [26] & tilt test; geometry unclear; recorded as M03 static upper\\
glazed ceramic on melamine-faced base & 0.28 & static & [26] & tilt test; recorded as M08 static upper\\
cardboard (untreated) on wooden pallet & 0.50 & static-derived design value & [27] & lashing-guide factor; recorded as M14 static upper\\
unpainted rough steel on sawn wood & 0.45 & static-derived design value & [27] & lashing-guide factor; recorded as M18 static upper\\
oak on oak, grain along (static) & 0.62 & static & [21] & secondhand; recorded as V02 static upper\\
oak across on oak along (static) & 0.54 perp. / 0.43 crossed & static & [21] & secondhand; V03 static upper follows the Section 10 decision\\
PLEXIGLAS (cast acrylic), datasheet & 0.8 and 0.5 & unlabelled & [28] & counterface and static/kinetic not stated; cannot be used (M06, M15, M24 stay none)\\
wear-plate steel on bare (unpainted) MDF & 0.68 & kinetic, excluded & [17] & bare MDF; paint changes the surface, so not used for M27\\
polystyrene on steel & 0.30--0.35 & static & [3] & v1 excluded list\\
Plexiglas (acrylic) on steel & 0.40--0.50 & static & [3] & v1 excluded list\\
nylon on steel & 0.40 & static & [3] & v1 excluded list\\
polyethylene on steel & 0.20 & static & [3] & v1 excluded list\\
wood on clean metal & 0.20--0.60 & static & [3] & v1 excluded list (the range the 8 Oct draft used as kinetic)\\
wood on wood & 0.25--0.50 & static & [3] & v1 excluded list\\
leather on wood & 0.30--0.40 & static & [3] & v1 excluded list\\
paper on cast iron & 0.20 & static & [3] & v1 excluded list\\
PLA (3D printed) on steel & approx. 0.30 static & static & [9] & v1 excluded list, derived from reported forces; no kinetic value\\

\end{longtable}
}

\section{Table 3: rig pairs}\label{sec:rig}
One row per rig pair, with the pair IDs of the measured file: $\mathrm{M} = 9(\text{surface}-1) + \text{object}$, so M18 is O9 on S2. Objects O1--O9: PLA (3D printed), ABS sheet (extruded), furniture felt, neoprene rubber 60A, corrugated cardboard, cast acrylic, oak (grain along travel), glazed ceramic tile, stainless steel 304; O7b is oak with grain across travel. Surfaces S1--S5: white melamine shelf board, birch plywood (sanded), painted MDF (satin), oak board (grain along travel), ABS sheet (smooth face). Neoprene is expected to exceed 0.5, so it will probably fall outside the simulated range (0.10--0.50) regardless of the entry.

{\footnotesize
\begin{longtable}{l p{2.3cm} p{2.4cm} r c c c l p{3.9cm}}
\caption{Rig table. Tier none means no kinetic source: the value columns are blank in the CSV and C2 falls back to C0. ``Static upper'' is a nearby static value, never an entry.}\label{tab:rig}\\
\hline
ID & Object & Surface & $\mu_k$ & Tier & 90\% interval & Static upper & Src & Note\\
\hline
\endfirsthead
\multicolumn{9}{l}{\emph{Table~\ref{tab:rig}, continued}}\\
\hline
ID & Object & Surface & $\mu_k$ & Tier & 90\% interval & Static upper & Src & Note\\
\hline
\endhead
\hline
\endfoot
M01 & PLA (3D printed) & white melamine shelf board & 0.29 & C & 0.12--0.73 &  & [11] & PLA on acrylic, inclined-plane launch\\
M02 & ABS sheet (extruded) & white melamine shelf board & 0.30 & C & 0.12--0.75 &  & [13,14,12] & Printed ABS on steel, 0.05--0.75 m/s (sheet ABS on wood/laminate not found: different counterface class)\\
M03 & furniture felt & white melamine shelf board & --- & none & --- & 0.41 & [26] & Felt and ceramic on a melamine base; geometry unclear; static only\\
M04 & neoprene rubber 60A & white melamine shelf board & 0.83 & C & 0.33--2.08 &  & [22] & Rubber sole on varnished maple; rubber friction rises with speed [23]; above the simulated range\\
M05 & corrugated cardboard & white melamine shelf board & 0.20 & C & 0.08--0.50 &  & [24] & Cardboard on resin-coated board, sliding\\
M06 & cast acrylic & white melamine shelf board & --- & none & --- &  & --- & No kinetic source\\
M07 & oak, grain along travel & white melamine shelf board & 0.39 & C & 0.16--0.98 &  & [20] & Dry wood on hard smooth surfaces, 0.3--0.5 (geometric mean)\\
M08 & glazed ceramic tile & white melamine shelf board & --- & none & --- & 0.28 & [26] & Ceramic on melamine base, tilt test; static only\\
M09 & stainless steel 304 & white melamine shelf board & 0.25 & B & 0.125--0.50 &  & [17] & Wear-plate steel on decor-faced chipboard, 22 mm/s\\
M10 & PLA (3D printed) & birch plywood, sanded & 0.33 & B & 0.165--0.66 &  & [11] & PLA on fine-sanded wood\\
M11 & ABS sheet (extruded) & birch plywood, sanded & 0.30 & C & 0.12--0.75 &  & [13,14,12] & As M02\\
M12 & furniture felt & birch plywood, sanded & --- & none & --- &  & --- & No source\\
M13 & neoprene rubber 60A & birch plywood, sanded & 0.83 & C & 0.33--2.08 &  & [22] & As M04\\
M14 & corrugated cardboard & birch plywood, sanded & 0.21 & B & 0.105--0.42 & 0.50 & [25,24,27] & Corrugated box on smooth pine; static upper from CTU design value\\
M15 & cast acrylic & birch plywood, sanded & --- & none & --- &  & --- & No kinetic source\\
M16 & oak, grain along travel & birch plywood, sanded & 0.34 & C & 0.14--0.85 &  & [18,21] & Geometric mean of timber on plywood (0.24) and oak on oak (0.48)\\
M17 & glazed ceramic tile & birch plywood, sanded & --- & none & --- &  & --- & No source\\
M18 & stainless steel 304 & birch plywood, sanded & 0.22 & C & 0.088--0.55 & 0.45 & [18,19,27] & Smooth steel and timber, reversed roles; static upper from CTU design value\\
M19 & PLA (3D printed) & painted MDF (satin) & 0.29 & C & 0.12--0.73 &  & [11] & PLA on acrylic\\
M20 & ABS sheet (extruded) & painted MDF (satin) & 0.30 & C & 0.12--0.75 &  & [13,14,12] & As M02\\
M21 & furniture felt & painted MDF (satin) & --- & none & --- &  & --- & No source\\
M22 & neoprene rubber 60A & painted MDF (satin) & 0.83 & C & 0.33--2.08 &  & [22] & Varnished surface; as M04\\
M23 & corrugated cardboard & painted MDF (satin) & 0.20 & C & 0.08--0.50 &  & [24] & Coated board; as M05\\
M24 & cast acrylic & painted MDF (satin) & --- & none & --- &  & --- & No kinetic source\\
M25 & oak, grain along travel & painted MDF (satin) & 0.39 & C & 0.16--0.98 &  & [20] & As M07\\
M26 & glazed ceramic tile & painted MDF (satin) & --- & none & --- &  & --- & No source\\
M27 & stainless steel 304 & painted MDF (satin) & --- & none & --- &  & --- & Bare-MDF value (0.68 [17]) excluded: paint changes the surface\\
V01 & stainless steel 304 & ABS sheet, smooth face & 0.24 & C & 0.096--0.60 &  & [12] & Same source as Y01 (L21)\\
V02 & oak, grain along travel & oak board, grain along travel & 0.48 & A & 0.32--0.72 & 0.62 & [21] & Morin and Kollmann values, secondhand; same figure as handbook L17\\
V03 & oak, grain across travel & oak board, grain along travel & 0.19 & A & 0.13--0.29 & 0.43 & [21] & PLACEHOLDER (Section 10 decision 2): 'crossed' 0.19 (90\% 0.13--0.29, static 0.43) used; 'perpendicular' 0.33 (90\% 0.22--0.50, static 0.54) is the alternative; team decides\\

\end{longtable}
}

\begin{table}[htbp]
\centering\small
\caption{Rig coverage at a glance: $\mu_k$ (tier) per object and surface, built from the M rows; ``---'' = tier none, with any static upper bound in parentheses. 17 of 27 cells are covered (B: 3, C: 14, A: 0).}\label{tab:grid}
\begin{tabular}{l l l l}
\hline
Object & S1 melamine & S2 birch plywood & S3 painted MDF\\
\hline
O1 PLA (3D printed) & 0.29 (C) & 0.33 (B) & 0.29 (C)\\
O2 ABS sheet (extruded) & 0.30 (C) & 0.30 (C) & 0.30 (C)\\
O3 furniture felt & --- (static 0.41) & --- & ---\\
O4 neoprene rubber 60A & 0.83 (C) & 0.83 (C) & 0.83 (C)\\
O5 corrugated cardboard & 0.20 (C) & 0.21 (B) & 0.20 (C)\\
O6 cast acrylic & --- & --- & ---\\
O7 oak, grain along travel & 0.39 (C) & 0.34 (C) & 0.39 (C)\\
O8 glazed ceramic tile & --- (static 0.28) & --- & ---\\
O9 stainless steel 304 & 0.25 (B) & 0.22 (C) & ---\\

\hline
\end{tabular}
\end{table}

\section{Separating Yu from the lookup}\label{sec:yu}
Yu et al.\ (2016) set the simulator friction for Y01--Y03 (Y04 was dropped on 5 Nov, decision 6), so their numbers cannot also be lookup entries (rule 4). In v1 they were rows L21--L24; in v2 the same pair IDs as the truth file (\texttt{ground\_truth\_published.csv}) get independent values from the sources below, and the Yu figures are kept out of every row and source list of \texttt{friction\_lookup.csv}. Table~\ref{tab:yu} is the only place in the lookup documents where the two are side by side; its ``Yu truth'' column is read from the truth file when the row file is generated.

\begin{table}[htbp]
\centering\footnotesize
\caption{Yu pairs: simulator truth vs.\ the independent lookup value. Ratio = lookup / truth.}\label{tab:yu}
\begin{tabular}{p{4.0cm} c c c c c p{4.6cm}}
\hline
Pair & Yu truth (simulator only) & Independent lookup & Tier & 90\% interval & Ratio & Source and match\\
\hline
Y01 (L21) stainless steel block on ABS sheet & 0.14 & 0.24 & C & 0.096--0.60 & 1.7$\times$ & [12]: printed ABS under a stainless hemisphere, 0.05 m/s\\
Y02 (L22) stainless steel block on acetal (Delrin) sheet & 0.14 & 0.31 & A & 0.21--0.47 & 2.2$\times$ & [15]: Delrin block on steel ring, 0.5 m/s, 10 N\\
Y03 (L23) stainless steel block on plywood sheet & 0.26 & 0.22 & C & 0.088--0.55 & 0.85$\times$ & [18], [19]: smooth steel and timber, reversed roles\\
Y04 (L24) stainless steel block on polyurethane sheet & 0.285 & 0.61 & B & 0.31--1.22 & 2.1$\times$ & [16]: TPU 95A on brushed steel, flat on flat, 0.32 m/s\\

\hline
\end{tabular}
\end{table}

\noindent Y02 (L22) is the clearest case: exact materials at a matching speed, still about $2\times$ off ([15] alone gives 0.31; the pooled entry of decision 5, 0.25, is $1.8\times$). The datasheet value [15] comes from a loaded block-on-ring test, while Yu slid a lightly loaded flat block. That is the ASTM G115 point: a friction coefficient belongs to the whole test system, not to a pair of material names.

\medskip\noindent In row L06, \texttt{yu} is removed from \texttt{sources} and the Yu figure from \texttt{notes}; [15] (0.30--0.33 at 0.5~m/s) is added as a further source, and L06 and Y02 share one pooled acetal-on-steel entry, 0.25 with $f = 1.8$ (decision 5, Section~\ref{sec:placeholders}); the v1 nominal 0.20 (ASTM D3702) would have been kept had the choice been binary.

\medskip\noindent Rows L01--L20 stay. Their simulator truth is sampled from the table's own interval, so C2 wins there by construction. Report them as a memorisation probe, not as evidence that C2 is accurate.

\section{Difficulties (v2 list, to be merged with v1)}\label{sec:difficulties}
\begin{enumerate}
 \item \textbf{Kinetic vs static.} Most compilations do not say which they give, and static is usually the larger. Tilt tests (Papalou [26]) and lashing-guide factors (CTU Code [27]) are static by construction. Rule 1 keeps them out of the entries; \texttt{static\_upper} keeps the information.
 \item \textbf{Speed dependence.} Rubber-like materials get stickier with speed (Roth et al.\ [23]); the TPU value behind the dropped Y04 is at 0.32~m/s and the published PLA-on-steel values span 0.23 to 0.49 across speeds, which is where spread B comes from. The CSV records the source speed where the paper states one; it is blank for handbooks, inclined planes and the unreadable Valiant scan.
 \item \textbf{A friction coefficient belongs to the whole test system (ASTM G115).} Y02 (L22) is the example: Delrin on steel at 0.5~m/s from a loaded block-on-ring test [15] gives 0.31, Yu's lightly loaded flat block gave about half that. Same materials, same speed, different system, $2\times$ apart.
 \item \textbf{Extraction-garbled tables.} The Meng 2008 Table~3 rows behind M16, M18 and Y03 came out of text extraction with uncertain row assignment, and the Valiant 1986 scan is too poor to read the speed. Both need the PDF opened (Section~\ref{sec:verify}).
 \item \textbf{Secondhand values.} The oak-on-oak numbers (V02, V03, M16) are Kollmann and Morin values as tabulated by Aira 2014 and Villar-Garc\'{\i}a 2020; Kollmann 1959 itself was not opened. Mens 1991 reached us via a Curbell white paper and the NASA table via engineeringlibrary.org.
 \item \textbf{The 0.10--0.50 simulated range.} Neoprene (0.83) and TPU (0.61, the dropped Y04) are above it, so for those pairs C2 is wrong in the simulator whatever the source says; the simulated pairs are chosen inside the range and the lookup is not.
 \item \textbf{Circularity, L01--L20 only.} For the handbook rows the simulator truth is sampled from the table's own 90\% interval, so C2 is correct there by construction. Report those pairs as a memorisation probe. Y01--Y03 and the rig rows are no longer circular: their truth comes from Yu or from our measurements, and their lookup values from other sources.
 \item \textbf{Tier spreads are a modelling choice.} The factors 2, 2.5 and 3 (decision 3, 5 Nov) are not fitted to anything beyond the PLA-on-steel spread.
 \item \textbf{Gaps.} Felt, cast acrylic and glazed tile have no kinetic source on any rig surface; stainless on painted MDF has none either because the only MDF value is for bare MDF. Ten of 27 rig cells are tier none.
\end{enumerate}

\section{Verify before citing}\label{sec:verify}
Set \texttt{verified = y} in \texttt{lookup\_table.py} only after checking these against the PDF:
\begin{itemize}
 \item Meng 2008 Table~3 [18]: row assignments for M16, M18 and Y03 (L23) were garbled in extraction.
 \item Kukla 2021 [17]: that ``laminated chipboard'' is melamine-faced (M09).
 \item Bories Otondo 2024 [11]: which print face touched the plane (M01, M10, M19).
 \item Oak-on-oak values [21]: secondhand via Aira 2014 and Villar-Garc\'{\i}a 2020; Kollmann 1959 not opened (V02, V03, M16).
 \item DuPont L-12259 [15]: test geometry and footnotes (Y02 / L22, and the L06 addition).
 \item Mahmood 2023 [12]: published values vs the preprint (Y01 / L21, V01, M02).
 \item Valiant 1986 [22]: poor scan, speed unreadable (M04, M13, M22).
 \item GDV Container Handbook [24]: Tables 1c--1d (M05, M14, M23).
 \item Murase 1984 (figures only, not yet a source): Figs.~7--13 could fill the acrylic and stainless rows.
 \item Reference titles marked ``[title: confirm]'' in Section~\ref{sec:refs} were not recorded when the values were taken; the PLEXIGLAS datasheet (Table~\ref{tab:static}) has no URL on file.
\end{itemize}

\section{Team decisions (spec Section 10, settled 5 Nov 2026)}\label{sec:placeholders}
Decided by the team on 5 Nov 2026 and applied in \texttt{lookup\_table.py}, which prints them on every run.
\begin{itemize}
 \item \textbf{Decision 2, V03 (oak across on oak along).} Nominal = Kollmann ``perpendicular'' 0.33 (source 0.32--0.34; 90\% 0.165--0.66 at $f = 2$; static upper 0.54). The ``crossed'' 0.19 (static 0.43) is kept in the row's notes as the alternative.
 \item \textbf{Decision 3, tier spread factors.} $f = 2 / 2.5 / 3$ for A / B / C, entered as one observation of $\ln\mu$ on top of the population prior (\texttt{belief.loglik\_lookup}).
 \item \textbf{Decision 5, L06 and Y02.} One pooled acetal-on-steel entry, 0.25 with $f = 1.8$ (90\% 0.14--0.45), from [4], [5], [7] and [15]; 0.20 (ASTM D3702) would have been kept had the choice been binary.
 \item \textbf{Decision 6, Y04 (polyurethane).} Dropped from the simulated set: the row is gone from \texttt{ground\_truth\_published.csv} and the lookup candidate (0.61 [16], above the cap) is in Table~\ref{tab:excluded}.
\end{itemize}

\section{Other excluded candidates}
Candidates from v1 dropped for reasons other than being static: above the 0.50 cap or below the 0.10 floor, toxic or exotic, or a range too wide to use.

{\footnotesize
\begin{longtable}{p{4.8cm} l l p{8cm}}
\caption{Excluded candidates (non-static reasons).}\label{tab:excluded}\\
\hline
Pair & Value & Src & Reason\\
\hline
\endfirsthead
\multicolumn{4}{l}{\emph{Table~\ref{tab:excluded}, continued}}\\
\hline
Pair & Value & Src & Reason\\
\hline
\endhead
\hline
\endfoot
stainless steel on polyurethane (TPU 95A) sheet, 0.32 m/s & 0.61 & [16] & above 0.50 cap; Yu pair Y04, dropped from the simulated set on 5 Nov (decision 6); rubber-like friction rises with speed\\
mild steel on mild steel & 0.57 & [1] & above 0.50 cap\\
nickel on nickel & 0.53 & [1] & above 0.50 cap\\
nickel on mild steel & 0.64 & [1] & above 0.50 cap\\
glass on nickel & 0.56 & [1] & above 0.50 cap\\
copper on glass & 0.53 & [1] & above 0.50 cap\\
leather on oak (parallel) & 0.52 & [1] & above 0.50 cap\\
leather on cast iron & 0.56 & [1] & above 0.50 cap\\
aluminium on aluminium & 1.40 & [1] & above cap; galling\\
mild steel on lead & 0.95 & [1] & above cap\\
rubber on dry concrete / asphalt & 0.5--0.85 & [3] & above cap; range only\\
PTFE on PTFE / on steel & 0.04 & [1] & below 0.10 floor; slides too long\\
waxed wood on snow, ice on ice & 0.02--0.10 & [3] & below floor; not a table surface\\
lead on cast iron & 0.43 & [1] & toxic; excluded for realism (value is usable)\\
tin on cast iron & 0.32 & [1] & exotic; dropped to reach 24 (value is usable)\\
cadmium on mild steel & 0.46 & [1] & toxic; excluded\\
magnesium on mild steel & 0.42 & [1] & dropped to reach 24 (value is usable)\\
hard steel on babbitt & 0.33--0.35 & [1] & bearing alloy; exotic\\
mild steel on phosphor bronze & 0.34 & [1] & dropped to reach 24 (value is usable)\\
garnet on mild steel & 0.39 & [1] & abrasive mineral; not a plausible object\\
ABS (generic) & 0.3--1.0 & [6] & range too wide to use; static/dynamic unspecified\\
acrylic, cast (generic) & 0.04--0.5 & [6] & range too wide to use; static/dynamic unspecified\\
HDPE (generic) & 0.22--0.62 & [6] & range too wide; static/dynamic unspecified\\
paper / board on paper & 0.2--0.9 & [10] & paper-on-paper only; kinetic only in figures\\

\end{longtable}
}

\section{References}\label{sec:refs}
Numbers are those used in the tables; the bold key is the \texttt{sources} token in \texttt{friction\_lookup.csv}. Sources 1--10 are the v1 keys (9 and 10 appear only in the excluded lists); 11--27 are the v2 lookup sources; 28 appears only in Table~\ref{tab:static}.
\begingroup
\renewcommand{\section}[2]{}% swallow thebibliography's own \section*{\refname}; the heading above is used instead
\begin{thebibliography}{99}
\raggedright
\bibitem[1]{nasa} \textbf{nasa} --- NASA RP-1228, R. T. Barrett, \emph{Fastener Design Manual}, 1990, Table IV (coefficients of friction, dry sliding). \url{https://ntrs.nasa.gov/citations/19900009424}
\bibitem[2]{elib} \textbf{elib} --- engineeringlibrary.org, reproduction of NASA RP-1228 Table IV (used to read the table; same numbers). \url{https://engineeringlibrary.org/}
\bibitem[3]{etb} \textbf{etb} --- The Engineering ToolBox, \emph{Friction and Friction Coefficients} (uncited compilation; cross-check only, rule 5). \url{https://www.engineeringtoolbox.com/friction-coefficients-d_778.html}
\bibitem[4]{delrin} \textbf{delrin} --- DuPont, \emph{Delrin 150 acetal homopolymer} datasheet: coefficient of friction by ASTM D3702 (thrust washer, 40 psi, 50 fpm $\approx$ 0.25 m/s).
\bibitem[5]{tse} \textbf{tse} --- TSE Industries, \emph{Physical Properties of Plastics} (table of dynamic coefficient of friction against steel).
\bibitem[6]{mcm} \textbf{mcm} --- McMaster-Carr, \emph{More About Plastics}, 2016 (friction ranges, static/dynamic not always stated).
\bibitem[7]{mens} \textbf{mens} --- J. W. M. Mens and A. W. J. de Gee, \emph{Friction and wear behaviour of 18 polymers in contact with steel in environments of air and water}, Wear 149 (1991) 255--268; values via a Curbell Plastics white paper.
\bibitem[8]{villar} \textbf{villar} --- J. R. Villar-Garc\'{\i}a et al., chestnut timber friction by direct shear, Int. Agrophysics 34, 2020 [title: confirm].
\bibitem[9]{stoimenov} \textbf{stoimenov} --- N. Stoimenov et al., friction of 3D-printed PLA on steel; static value derived from reported forces [full citation: confirm]. Excluded-list only.
\bibitem[10]{silveira} \textbf{silveira} --- Silveira et al., paper and board friction (paper on paper; kinetic only in figures) [full citation: confirm]. Excluded-list only.
\bibitem[11]{BO} \textbf{BO} --- Bories Otondo, Montes de Oca and Piovan, 3D-printed PLA friction by inclined-plane launch, AJEA (UTN) 2024 [title: confirm]. \url{https://rtyc.utn.edu.ar/index.php/ajea/article/download/1851/1685/8895}
\bibitem[12]{MA} \textbf{MA} --- Mahmood et al., tribology of 3D-printed ABS against a stainless hemisphere, Tribology Online 18(6), 2023 [title: confirm]; preprint \url{https://www.researchsquare.com/article/rs-2281753/v1}
\bibitem[13]{MI} \textbf{MI} --- Mileti\'c et al., tribological behaviour of 3D-printed ABS on steel, SERBIATRIB '25 [title: confirm]. \url{https://scidar.kg.ac.rs/bitstream/123456789/22509/1/184.pdf}
\bibitem[14]{DJ} \textbf{DJ} --- Djurovi\'c et al., \emph{Tribological Behaviour and Surface Roughness Quality of 3D Printed ABS Material}, SERBIATRIB '23. \url{https://scidar.kg.ac.rs/bitstream/123456789/21541/1/Tribological%20Behaviour%20and%20Surface%20Roughness%20Quality%20of%203D%20Printed%20ABS%20Material.pdf}
\bibitem[15]{DU} \textbf{DU} --- DuPont, \emph{Delrin low wear / low friction grades}, brochure L-12259 (block-on-ring data against steel). \url{https://www.emcoplastics.com/assets/pdf/delrin/Delrin%20Low%20Wear%20Low%20Friction%20Brochure%20Dupont.pdf}
\bibitem[16]{MZ} \textbf{MZ} --- Mart\'{\i}nez et al., TPU--steel long-stroke sliding tests, LubMat '12 [title: confirm]. \url{https://core.ac.uk/download/16754648.pdf}
\bibitem[17]{KU} \textbf{KU} --- Kukla, Wargu\l{}a and Biszczanik, \emph{Determining the Coefficient of Friction of Wood-Based Materials for Furniture Panels in the Aspect of Modelling Their Shredding Process}, Wood Research 66(5), 2021. \url{https://www.woodresearch.sk/cms/wp-content/uploads/2021/10/Determining-The-Coefficient-Of-Friction-Of-Wood-Based-Materals-For-Furniture-Panels-In-The-Aspect-Of-Modelling-Their-Shredding-Process.pdf}
\bibitem[18]{ME} \textbf{ME} --- Meng, Hirai and Koizumi, friction between timber and steel (Table 3), Mokuzai Gakkaishi 54(5), 2008 [title: confirm]. \url{https://eprints.lib.hokudai.ac.jp/repo/huscap/all/68014/54_281.pdf}
\bibitem[19]{VI} \textbf{VI} --- Villar-Garc\'{\i}a et al., timber friction coefficients, Forests 13(7), 1078, 2022 [title: confirm]. \url{https://www.mdpi.com/1999-4907/13/7/1078}
\bibitem[20]{WH} \textbf{WH} --- S. V. Glass and S. L. Zelinka, \emph{Moisture Relations and Physical Properties of Wood}, Ch. 4 in \emph{Wood Handbook: Wood as an Engineering Material}, FPL-GTR-190, 2010, p. 4-17 (friction). \url{https://archive.org/details/wood-as-an-engineering-material}
\bibitem[21]{KO} \textbf{KO} --- Kollmann (1959) and Morin values for wood-on-wood friction, as tabulated secondhand in Aira et al., Mater. Construcc. 64(315), 2014 (\url{https://materconstrucc.revistas.csic.es/index.php/materconstrucc/article/download/1482/1746?inline=1}) and Villar-Garc\'{\i}a et al., Int. Agrophysics 34, 2020. Kollmann not opened.
\bibitem[22]{VA} \textbf{VA} --- Valiant, friction of shoe-sole rubber on sports floors (varnished maple), ISBS 1986 [title: confirm; poor scan, speed unreadable]. \url{https://ojs.ub.uni-konstanz.de/cpa/article/view/1497/1347}
\bibitem[23]{RO} \textbf{RO} --- F. L. Roth, R. L. Driscoll and W. L. Holt, \emph{Frictional properties of rubber}, J. Res. NBS 28 (1942) 439. \url{https://nvlpubs.nist.gov/nistpubs/jres/28/jresv28n4p439_A1b.pdf}
\bibitem[24]{GDV} \textbf{GDV} --- GDV (German Insurance Association), \emph{Container Handbook}, Sec. 4.4.5, friction tables 1c--1d. \url{https://www.containerhandbuch.de/chb_e/stra/stra_04_04_05.html}
\bibitem[25]{YA} \textbf{YA} --- Yang, Luo and Cai, friction of corrugated boxes on wood, WCMNM 2017 [title: confirm]. \url{https://core.ac.uk/download/82919615.pdf}
\bibitem[26]{PA} \textbf{PA} --- Papalou, tilt-test friction coefficients of common materials, Applied Mechanics 5(4), 48, 2024 [title: confirm]. \url{https://www.mdpi.com/2673-3161/5/4/48}
\bibitem[27]{CTU} \textbf{CTU} --- IMO/ILO/UNECE \emph{CTU Code}, Quick Lashing Guide (IMO MSC.1/Circ.1498) and the Cordstrap CTU Code guide: static-derived design friction factors (not kinetic measurements).
\bibitem[28]{PLX} \textbf{PLX} --- R\"{o}hm, PLEXIGLAS (cast acrylic) product data: friction values 0.8 and 0.5 quoted without counterface or static/kinetic label [datasheet and URL: locate before citing]. Table 2 only.

\end{thebibliography}
\endgroup

\medskip\noindent\textbf{Not a lookup source (rule 4).} K.-T. Yu, M. Bauza, N. Fazeli and A. Rodriguez, \emph{More than a million ways to be pushed: a high-fidelity experimental dataset of planar pushing}, IROS 2016. Sets the simulator friction for Y01--Y03 via \texttt{ground\_truth\_published.csv}; cited here only to explain Table~\ref{tab:yu}.

\end{document}
